%% Plantilla elsarticle — Expert Systems with Applications (Q1)
%% Artículo: Framework Multi-Agente LLM para Automatización de Procurement
%% Autora: Edurne Martínez de Contrasta | PROJENER.AI SL · UAX 2026

\documentclass[preprint,12pt,authoryear]{elsarticle}

%% ── Paquetes ─────────────────────────────────────────────────────────────────
\usepackage[utf8]{inputenc}
\usepackage[T1]{fontenc}
\usepackage[english]{babel}
\usepackage{amsmath,amssymb}
\usepackage{graphicx}
\usepackage{booktabs}
\usepackage{multirow}
\usepackage{array}
\usepackage{longtable}
\usepackage{tabularx}
\usepackage{xcolor}
\usepackage{hyperref}
\usepackage{url}
\usepackage{siunitx}
\usepackage{caption}
\usepackage{subcaption}
\usepackage{float}
\usepackage{enumitem}
\usepackage{listings}
\usepackage{lscape}

\lstset{
  basicstyle=\ttfamily\small,
  breaklines=true,
  frame=single,
  numbers=left,
  numberstyle=\tiny,
  keywordstyle=\color{blue},
  stringstyle=\color{red!60!black},
  commentstyle=\color{gray},
  showstringspaces=false
}

\journal{Expert Systems with Applications}

\begin{document}

\begin{frontmatter}

\title{Un Framework Multi-Agente LLM para la Automatizaci\'on Aut\'onoma
       de Procesos Empresariales: Caso de Estudio en Procurement}

\author[uax,proj]{Edurne Mart\'inez de Contrasta}
\ead{edurne.martinez@uax.es}

\address[uax]{M\'aster en Big Data y Business Analytics,
              Universidad Alfonso X el Sabio, Madrid, Espa\~na}
\address[proj]{PROJENER.AI SL, Madrid, Espa\~na}

\begin{abstract}
Los sistemas tradicionales de Automatizaci\'on Rob\'otica de Procesos (RPA)
automatizan procesos empresariales mediante reglas if-then, pero carecen de
capacidad de razonamiento ante excepciones.
Este trabajo propone, implementa y eval\'ua un framework multi-agente basado en
Modelos de Lenguaje de Gran Escala (LLMs) para la automatizaci\'on end-to-end
de procesos de procurement empresarial.
El sistema implementa cinco agentes especializados con roles empresariales
(\emph{Requester}, \emph{Procurement}, \emph{Finance}, \emph{Legal},
\emph{Compliance}) orquestados con CrewAI, cada uno con acceso exclusivo a sus
propias herramientas a trav\'es de seis APIs mock que simulan sistemas ERP,
plataformas CRM y bases de datos regulatorias.

El experimento compara ocho variantes del sistema (M1--M5, M2x, M3b, M4b) sobre
un conjunto de 50 casos de procurement en tres niveles de complejidad, utilizando
los modelos \texttt{llama-3.1-8b-instant} y \texttt{llama-3.3-70b-versatile} de
Groq. Los resultados revelan cuatro hallazgos principales:
\begin{enumerate}[label=(\arabic*)]
  \item el agente \'unico con contexto expandido (M2x: ARR\,=\,84\%) supera al
        sistema multi-agente hom\'ogneo con igual presupuesto de tokens
        (M3: ARR\,=\,68\%), validando emp\'iricamente la hip\'otesis de
        \emph{compute equivalence}; sin embargo, el an\'alisis de robustez
        muestra alta variabilidad en M2 (ARR\,=\,92\%\,$\pm$\,11.3\,pp)
        frente a la estabilidad de M4 (ARR\,=\,98.7\%\,$\pm$\,0.9\,pp);
  \item introducir diversidad de modelos produce el mayor salto de rendimiento
        (+32\,pp ARR, de M3 a M4);
  \item el ablation study demuestra que eliminar los agentes Legal y Procurement
        aumenta el ARR en 20\,pp pero anula la HDA, evidenciando el trade-off
        entre automatizaci\'on y detecci\'on de riesgo; y
  \item la reproducibilidad se confirma al 100\% con \texttt{temperature=0.0}.
\end{enumerate}
El hallazgo central es que la arquitectura multi-agente requiere modelos con
capacidad suficiente para beneficiarse de la especializaci\'on.
\end{abstract}

\begin{keyword}
Sistemas multi-agente \sep
Modelos de Lenguaje de Gran Escala \sep
Automatizaci\'on de procurement \sep
CrewAI \sep
Human-in-the-Loop \sep
Compute equivalence \sep
Diversidad de modelos \sep
IA empresarial \sep
Ablation study \sep
Reproducibilidad
\end{keyword}

\end{frontmatter}

%% ─────────────────────────────────────────────────────────────────────────────
\section{Introducci\'on}
\label{sec:introduccion}

Cons\'iderese el siguiente escenario: un t\'ecnico de mantenimiento necesita
adquirir una pieza de repuesto valorada en 150 euros.
El proveedor est\'a homologado, el presupuesto disponible existe y la necesidad
es inmediata.
A pesar de ello, la solicitud tardar\'a entre tres y siete d\'ias laborables en
recibir una decisi\'on, recorriendo secuencialmente los departamentos de compras,
legal, finanzas y compliance.
El coste medio de tramitar ese pedido asciende a 112 euros adicionales en tiempo
de personal \citep{hackett2023}.
El cuello de botella no reside en la burocracia como tal, sino en que los
sistemas dise\~nados para automatizarla no son capaces de razonar.

Los sistemas RPA actuales ejecutan con eficacia tareas estructuradas y
repetitivas, pero fracasan sistem\'aticamente ante cualquier desviaci\'on: un
proveedor con alerta de conflicto de inter\'es, un contrato con cl\'ausula de
exclusividad o un importe que supera el umbral de aprobaci\'on del \'area.

El campo de los sistemas multi-agente basados en LLMs ofrece una respuesta
t\'ecnica. Frameworks como AutoGen \citep{wu2023autogen}, CrewAI
\citep{crewai2024} y LangGraph permiten construir equipos de agentes que
colaboran, razonan y adaptan su comportamiento ante excepciones.
ChatDev \citep{qian2023chatdev} mostr\'o que equipos de agentes pueden superar
a agentes individuales en tareas complejas; \citet{park2023generative}
documentaron comportamientos sociales emergentes en comunidades de agentes LLM.

Este trabajo aborda la brecha de evidencia emp\'irica con cuatro contribuciones:
\begin{enumerate}[label=(\arabic*)]
  \item dise\~no e implementaci\'on de un framework multi-agente completo para
        procurement con cinco agentes especializados y seis herramientas
        empresariales simuladas;
  \item un dataset de 50 casos con ground truth anotado en tres niveles de
        complejidad, siguiendo la metodolog\'ia de WorkArena
        \citep{drouin2024workarena} y AgentBench \citep{liu2023agentbench};
  \item evaluaci\'on de ocho variantes del sistema, incluyendo ablation study
        ejecutado, comparativa de compute equivalence y an\'alisis de
        diversidad de modelos; y
  \item verificaci\'on de reproducibilidad y an\'alisis de robustez estad\'istica
        mediante tres r\'eplicas independientes.
\end{enumerate}

\paragraph{Pregunta de investigaci\'on.}
\textit{?`Puede un sistema multi-agente basado en LLMs automatizar un proceso
de procurement end-to-end, adaptarse a excepciones y superar a los enfoques RPA
tradicionales en calidad de decisi\'on --- y bajo qu\'e condiciones
arquitect\'onicas y de modelo se produce esa ventaja?}

El art\'iculo se estructura como sigue. La Secci\'on~\ref{sec:relacionado}
revisa el trabajo relacionado. La Secci\'on~\ref{sec:framework} describe la
arquitectura. La Secci\'on~\ref{sec:setup} detalla el setup experimental.
La Secci\'on~\ref{sec:resultados} presenta los resultados. Las
Secciones~\ref{sec:ablation}, \ref{sec:compute} y~\ref{sec:diversidad}
profundizan en ablation, compute equivalence y diversidad de modelos.
La Secci\'on~\ref{sec:discusion} discute los hallazgos y la
Secci\'on~\ref{sec:conclusiones} presenta las conclusiones.

%% ─────────────────────────────────────────────────────────────────────────────
\section{Trabajo Relacionado}
\label{sec:relacionado}

\subsection{Sistemas multi-agente LLM}

\citet{wu2023autogen} introdujeron AutoGen, framework de conversaci\'on
multi-agente que permite a los LLMs colaborar mediante mensajes estructurados.
CrewAI \citep{crewai2024} extiende este paradigma con roles, tareas y equipos
inspirados en organizaciones humanas.
\citet{park2023generative} demostraron comportamientos sociales emergentes en
comunidades de agentes LLM.
ChatDev \citep{qian2023chatdev} aplic\'o el paradigma al desarrollo de software,
mostrando que equipos de agentes especializados superan a agentes individuales.

\subsection{Automatizaci\'on de procesos empresariales}

Los sistemas RPA (UiPath, Automation Anywhere) automatizan tareas repetitivas
basadas en reglas pero fallan ante excepciones.
WorkArena \citep{drouin2024workarena} y TaskBench \citep{lu2023taskbench}
proporcionan entornos de evaluaci\'on para agentes en procesos empresariales.
ToolBench \citep{qin2023toolllm} demostr\'o que los LLMs pueden aprender a
usar miles de APIs reales.
APIBench \citep{patil2023gorilla} eval\'ua la capacidad de los LLMs para
generar llamadas a APIs correctas.
OSWorld \citep{xie2024osworld} eval\'ua agentes sobre sistemas operativos
completos; a diferencia de OSWorld, este trabajo se centra en procesos
empresariales estructurados con roles definidos.

\subsection{Compute equivalence y diversidad de modelos}

\citet{wu2023autogen} plantean que la comparaci\'on justa es ``5 agentes con
$X$ tokens vs.\ 1 agente con $5X$ tokens''.
Este trabajo ejecuta esa comparaci\'on emp\'iricamente.
En ensemble learning, la diversidad de modelos reduce la correlaci\'on de
errores \citep{dietterich2000ensemble}.

%% ─────────────────────────────────────────────────────────────────────────────
\section{Framework Propuesto}
\label{sec:framework}

\subsection{Arquitectura general}

El sistema se estructura como un Grafo Ac\'iclico Dirigido (DAG) donde cada
nodo representa un agente especializado y las aristas el flujo de informaci\'on.
El orquestador CrewAI gestiona la secuencia de ejecuci\'on.
En M5, el agente Compliance tiene reglas expl\'icitas de escalaci\'on HiL.

\begin{table}[htbp]
\centering
\caption{Endpoints de APIs mock, agentes propietarios y funci\'on en el sistema.}
\label{tab:apis}
\begin{tabularx}{\textwidth}{llX}
\toprule
\textbf{API} & \textbf{Endpoint} & \textbf{Funci\'on en el sistema} \\
\midrule
Cat\'alogo Proveedores &
  \texttt{GET /api/proveedores/\{id\}} &
  Valida estado y alertas de conflicto de inter\'es \\
Verificaci\'on Presupuestaria &
  \texttt{GET /api/presupuesto/\{dept\}} &
  Comprueba fondos y nivel de aprobaci\'on \\
Base Regulatoria &
  \texttt{GET /api/normativa/\{cat\}} &
  Detecta GDPR, CE, PRL, AESA, normativa sanitaria \\
Validaci\'on Contractual &
  \texttt{POST /api/contratos/validar} &
  Analiza cl\'ausulas de riesgo (exclusividad, IP, penalizaciones) \\
Formulario Solicitud &
  \texttt{POST /api/solicitudes/nueva} &
  Registra metadatos y abre el caso \\
Escalaci\'on HiL &
  \texttt{POST /api/escalacion/humano} &
  Escala con resumen, opciones y recomendaci\'on \\
\bottomrule
\end{tabularx}
\end{table}

\subsection{Las nueve variantes del sistema}

\begin{table}[htbp]
\centering
\caption{Resumen de variantes del sistema. M1, M2 y M4 reportan
         media\,$\pm$\,desv.\ t\'ipica (3 r\'eplicas). El resto reporta
         ejecuci\'on \'unica. ARR: Tasa de Resoluci\'on Aut\'onoma;
         HDA: Precisi\'on Detecci\'on HiL; DER: Tasa de Error; PT: Tiempo.}
\label{tab:variantes}
\resizebox{\textwidth}{!}{%
\begin{tabular}{llccccl}
\toprule
\textbf{Modelo} & \textbf{Config.} &
  \textbf{ARR} & \textbf{HDA} & \textbf{DER} &
  \textbf{PT (s)} & \textbf{Prop\'osito} \\
\midrule
M1 --- RPA & Reglas if-then &
  100\%\,$\pm$\,0.0\,pp & 0\%\,$\pm$\,0.0\,pp & 34\%\,$\pm$\,0.0\,pp &
  $\approx$0 & Baseline m\'inimo \\
M2 --- Single (200\,tok) & 8b\,$\times$\,1 ag. &
  92\%\,$\pm$\,11.3\,pp & 10.3\%\,$\pm$\,14.5\,pp & 27.3\%\,$\pm$\,3.8\,pp &
  0.31 & Referencia LLM \\
M2x --- Single (750\,tok) & 8b\,$\times$\,1 ag. &
  84\% & 23.1\% & 28\% & 0.40 & Compute equivalence \\
M3 --- Multi hom. & 8b\,$\times$\,5 ag. &
  68\% & 30.8\% & 28\% & 1.74 & Sistema principal \\
M3b --- Ablation 3\,ag & 8b\,$\times$\,3 ag. &
  88\% & 0\% & 30\% & 0.77 & Ablation Legal+Proc \\
M4 --- Multi mixto & 8b\,$\times$4+70b\,$\times$1 &
  98.7\%\,$\pm$\,0.9\,pp & 5.1\%\,$\pm$\,3.6\,pp & 27.3\%\,$\pm$\,0.9\,pp &
  1.69 & \textbf{\'OPTIMO} \\
M4b --- Multi potente & 70b\,$\times$\,5 ag. &
  98\% & 0\% & 26\% & 2.40 & Hom. potente \\
M5 --- Multi + HiL & 8b\,$\times$5 + HiL &
  98\% & 0\% & 28\% & 1.56 & HiL activo \\
M5b --- Clasif.\ Incert. & 2 llamadas\,$\times$\,1\,ag &
  94\% & 23.1\% & 22\% & 2.85 & Mejor HDA \\
\bottomrule
\end{tabular}%
}
\end{table}

%% ─────────────────────────────────────────────────────────────────────────────
\section{Setup Experimental}
\label{sec:setup}

\subsection{Dataset: 50 casos de procurement}

Se dise\~naron 50 casos con ground truth anotado en tres niveles:
\emph{Simple} (C01--C20), \emph{Medio} (C21--C38) y \emph{Complejo}
(C39--C50).
Siguiendo la metodolog\'ia de WorkArena \citep{drouin2024workarena} y
AgentBench \citep{liu2023agentbench}.

\subsection{M\'etricas}

\begin{description}[leftmargin=2cm, style=sameline]
  \item[ARR] Tasa de Resoluci\'on Aut\'onoma: \% casos resueltos sin HiL.
             Objetivo: $>$70\%.
  \item[HDA] Precisi\'on de Detecci\'on HiL: \% casos que requieren HiL y
             son escalados correctamente. Objetivo: $>$90\%.
  \item[DER] Tasa de Error de Decisi\'on: \% decisiones incorrectas vs.\
             ground truth. Objetivo: $<$10\%.
  \item[PT]  Tiempo de proceso medio por caso (segundos).
\end{description}

\subsection{Stack tecnol\'ogico}

\begin{itemize}
  \item Framework: CrewAI (Python 3.12, Windows 10, Jupyter Lab)
  \item LLM principal: \texttt{llama-3.1-8b-instant} v\'ia Groq API
  \item LLM decisor M4/M4b: \texttt{llama-3.3-70b-versatile} v\'ia Groq API
  \item Temperature: 0.0
  \item C\'odigo completo: \url{https://github.com/Edurne0781/multiagent-procurement-llm}
\end{itemize}

\subsection{An\'alisis de robustez estad\'istica}

Se ejecutaron tres r\'eplicas independientes de M1, M2 y M4.
M1 es perfectamente determin\'ista ($\sigma$\,=\,0).
M2 muestra alta variabilidad ($\sigma$\,ARR\,=\,11.3\,pp).
M4 es el m\'as estable ($\sigma$\,ARR\,=\,0.9\,pp).

\begin{figure}[htbp]
  \centering
  \includegraphics[width=\textwidth]{figura9}
  \caption{Robustez estad\'istica: media\,$\pm$\,desviaci\'on est\'andar de
           ARR, HDA y DER para M1, M2 y M4 sobre 3 r\'eplicas independientes.}
  \label{fig:replicas}
\end{figure}

\begin{table}[htbp]
\centering
\caption{Resultados individuales por r\'eplica para M1, M2 y M4.}
\label{tab:replicas}
\begin{tabular}{lcccc}
\toprule
\textbf{Modelo} & \textbf{R\'eplica 1} & \textbf{R\'eplica 2} &
  \textbf{R\'eplica 3} & \textbf{Media\,$\pm$\,$\sigma$} \\
\midrule
M1 --- ARR & 100\% & 100\% & 100\% & 100\%\,$\pm$\,0.0\,pp \\
M1 --- HDA & 0\%   & 0\%   & 0\%   & 0\%\,$\pm$\,0.0\,pp   \\
M1 --- DER & 34\%  & 34\%  & 34\%  & 34\%\,$\pm$\,0.0\,pp  \\
\midrule
M2 --- ARR & 76\%   & 100\% & 100\% & 92\%\,$\pm$\,11.3\,pp \\
M2 --- HDA & 30.8\% & 0\%   & 0\%   & 10.3\%\,$\pm$\,14.5\,pp \\
M2 --- DER & 22\%   & 28\%  & 32\%  & 27.3\%\,$\pm$\,3.8\,pp \\
\midrule
M4 --- ARR & 100\% & 98\%  & 98\%  & 98.7\%\,$\pm$\,0.9\,pp \\
M4 --- HDA & 0\%   & 7.7\% & 7.7\% & 5.1\%\,$\pm$\,3.6\,pp  \\
M4 --- DER & 26\%  & 28\%  & 28\%  & 27.3\%\,$\pm$\,0.9\,pp \\
\bottomrule
\end{tabular}
\end{table}

%% ─────────────────────────────────────────────────────────────────────────────
\section{Resultados}
\label{sec:resultados}

\subsection{Tabla comparativa global}

Ver Tabla~\ref{tab:variantes} para la comparaci\'on completa.
La Figura~\ref{fig:metricas} ofrece una visi\'on visual de las tres m\'etricas
principales.

\begin{figure}[htbp]
  \centering
  \includegraphics[width=\textwidth]{figura1}
  \caption{Comparativa de m\'etricas principales M1--M5: ARR (objetivo
           $>$70\%), HDA (objetivo $>$90\%) y DER (objetivo $<$10\%).}
  \label{fig:metricas}
\end{figure}

\subsection{An\'alisis por modelo}

\paragraph{M1 --- Baseline RPA.}
ARR\,=\,100\% porque nunca escala al humano.
HDA\,=\,0\% y DER\,=\,34\% (1 de cada 3 decisiones incorrecta).
Reproduce exactamente la limitaci\'on documentada de los sistemas RPA.

\paragraph{M2 --- Agente \'unico LLM (200 tokens).}
ARR\,=\,76\%, supera el objetivo $>$70\%.
DER\,=\,22\%, mejora de 12\,pp sobre M1.
HDA\,=\,30.8\%.
El modelo m\'as eficiente en tiempo (PT\,=\,0.26\,s).

\paragraph{M3 --- Multi-Agente hom\'ogneo (8b$\times$5).}
ARR\,=\,68\%, inferior a M2 (76\%).
Con \texttt{llama-3.1-8b-instant}, la fragmentaci\'on del contexto entre 5
agentes perjudica la calidad de decisi\'on.
La arquitectura multi-agente no aporta ventaja con modelos peque\~nos.

\paragraph{M4 --- Multi-Agente mixto (8b$\times$4 + 70b$\times$1).}
ARR\,=\,100\%, el mayor salto del experimento (+32\,pp sobre M3).
El agente Compliance con \texttt{llama-3.3-70b} toma decisiones
cualitativamente mejores.
Configuraci\'on \'optima en coste/rendimiento.

\paragraph{M4b --- Multi-Agente hom\'ogneo potente (70b$\times$5).}
ARR\,=\,98\%, pr\'acticamente igual que M4 pero 2.7$\times$ m\'as caro.
Confirma que los agentes auxiliares con 8b son suficientes.

\paragraph{M5 --- Multi-Agente + HiL activo.}
ARR\,=\,98\%. Las reglas expl\'icitas de escalaci\'on no aumentaron la HDA
(0\%), indicando que \texttt{llama-3.1-8b} carece de capacidad suficiente.

\paragraph{M5b --- Clasificador de incertidumbre.}
ARR\,=\,94\%, HDA\,=\,23.1\%, DER\,=\,22\%.
An\'alisis detallado en la Secci\'on~\ref{sec:ablation}.

\subsection{An\'alisis de coste computacional}

\begin{table}[htbp]
\centering
\caption{An\'alisis de coste computacional y tokens por modelo.}
\label{tab:tokens}
\begin{tabular}{lrrrrl}
\toprule
\textbf{Modelo} & \textbf{Tok/caso} & \textbf{Tok total} &
  \textbf{Llamadas} & \textbf{Coste equiv.} & \textbf{Observaci\'on} \\
\midrule
M1 --- RPA    &   0 &      0 & 0 & \$0.0000 & Sin LLM \\
M2 --- Single & 200 & 10.000 & 1 & \$0.0012 & M\'as eficiente \\
M2x           & 750 & 37.500 & 1 & \$0.0045 & Equiv. a M3 \\
M3 --- Multi  & 750 & 37.500 & 5 & \$0.0045 & M2x lo supera \\
M3b           & 450 & 22.500 & 3 & \$0.0027 & Ahorro 40\% \\
M4 --- Mixto  & 750 & 37.500 & 5 & \$0.0065 & ARR=100\% \\
M4b           & 750 & 37.500 & 5 & \$0.0120 & 2.7$\times$ M4 \\
M5 --- HiL    & 875 & 43.750 & 5 & \$0.0053 & Mejor ARR/coste \\
\bottomrule
\end{tabular}
\end{table}

\begin{figure}[htbp]
  \centering
  \includegraphics[width=\textwidth]{figura2}
  \caption{Curva de eficiencia: ARR frente al tiempo de proceso (PT) por caso.}
  \label{fig:eficiencia}
\end{figure}

\begin{figure}[htbp]
  \centering
  \includegraphics[width=\textwidth]{figura3}
  \caption{Distribuci\'on de decisiones por modelo sobre los 50 casos.}
  \label{fig:distribucion}
\end{figure}

\begin{figure}[htbp]
  \centering
  \includegraphics[width=\textwidth]{figura5}
  \caption{ARR frente al coste en tokens (izq.) y consumo de tokens por
           modelo (dcha.).}
  \label{fig:coste}
\end{figure}

%% ─────────────────────────────────────────────────────────────────────────────
\section{Ablation Study}
\label{sec:ablation}

\begin{table}[htbp]
\centering
\caption{Experimentos del ablation study.}
\label{tab:ablation}
\begin{tabularx}{\textwidth}{llllX}
\toprule
\textbf{ID} & \textbf{Experimento} & \textbf{ARR} & \textbf{HDA} &
  \textbf{Hallazgo principal} \\
\midrule
A1 & Sin Legal (planif.) & --- & --- &
  Casos con cl\'ausulas complejas pierden calidad \\
A2 & Sin Compliance (planif.) & --- & --- &
  Casos regulatorios no se bloquean \\
A3 & 3 ag vs 5 (M3b, ejec.) & 88\% vs 68\% & 0\% vs 30.8\% &
  +20\,pp ARR pero destruye HDA \\
A4 & Sin HiL (M3 vs M5) & 68\% vs 98\% & 30.8\% vs 0\% &
  HiL mejora ARR pero no HDA \\
\bottomrule
\end{tabularx}
\end{table}

\subsection{A3 ejecutado --- 3 agentes vs.\ 5 (M3b)}

Eliminar Legal y Procurement sube el ARR de 68\% a 88\% (+20\,pp) pero la HDA
cae de 30.8\% a 0\%.
Esto cuantifica el trade-off fundamental: la especializaci\'on no maximiza la
automatizaci\'on bruta sino la calidad supervisada de decisi\'on, exactamente
como predice la teor\'ia de ventaja comparativa de Ricardo
\citeyearpar{ricardo1817}.

\subsection{An\'alisis de grafos}

Finance es el nodo con mayor betweenness centrality (4), actuando como puente
entre la validaci\'on operativa y la decisi\'on regulatoria.
En t\'erminos de Granovetter \citeyearpar{granovetter1973}, Finance representa
el ``lazo d\'ebil'' m\'as cr\'itico del sistema.

\begin{table}[htbp]
\centering
\caption{M\'etricas de centralidad en la red de comunicaci\'on.}
\label{tab:grafos}
\begin{tabularx}{\textwidth}{lrrrX}
\toprule
\textbf{Agente} & \textbf{Grado total} & \textbf{Grado entrada} &
  \textbf{Betweenness} & \textbf{Rol} \\
\midrule
Requester   & 1 & 0 & 0 & Fuente \\
Procurement & 2 & 1 & 3 & Intermediario \\
Finance     & 2 & 1 & 4 & Nodo m\'as central \\
Legal       & 2 & 1 & 3 & Intermediario \\
Compliance  & 1 & 1 & 0 & Sumidero --- decide \\
\bottomrule
\end{tabularx}
\end{table}

\begin{figure}[htbp]
  \centering
  \includegraphics[width=\textwidth]{figura6}
  \caption{Red de comunicaci\'on entre agentes: topolog\'ia lineal (izq.) y
           centralidad por agente (dcha.).}
  \label{fig:grafo}
\end{figure}

\subsection{Convergencia por nivel de complejidad}

\begin{table}[htbp]
\centering
\caption{Tasa de resoluci\'on por nivel de complejidad.}
\label{tab:complejidad}
\begin{tabular}{lccccc}
\toprule
\textbf{Modelo} & \textbf{Simple} & \textbf{Medio} &
  \textbf{Complejo} & \textbf{Total} & \textbf{Observaci\'on} \\
\midrule
M2 Single    & 100\% & 100\%  & 100\%  & 100\% & Sin errores \\
M3 Multi     & 95\%  & 83.3\% & 83.3\% & 88\%  & Degradaci\'on en complejos \\
M4 Mixto     & 100\% & 100\%  & 100\%  & 100\% & Sin degradaci\'on \\
M5 HiL       & 100\% & 94.4\% & 100\%  & 98\%  & Un error rate-limit \\
M3b Ablation & 95\%  & 88.9\% & 83.3\% & 90\%  & M\'as robusto que M3 \\
\bottomrule
\end{tabular}
\end{table}

\begin{figure}[htbp]
  \centering
  \includegraphics[width=\textwidth]{figura7}
  \caption{Tasa de resoluci\'on por nivel de complejidad para cada modelo.}
  \label{fig:complejidad}
\end{figure}

\subsection{M5b --- Clasificador de incertidumbre}

M5b es un sistema de dos etapas: un agente Evaluador de Incertidumbre analiza
los factores de riesgo y emite un nivel de confianza (alta/baja); si es baja,
el caso se escala al humano; si es alta, un Gestor Senior decide.
HDA sube de 0\% a 23.1\% con 0 falsos positivos.

\begin{figure}[htbp]
  \centering
  \includegraphics[width=\textwidth]{figura8}
  \caption{M5 vs M5b: comparativa de m\'etricas (izq.) y an\'alisis de las
           13 escalaciones HiL requeridas (dcha.).}
  \label{fig:m5b}
\end{figure}

%% ─────────────────────────────────────────────────────────────────────────────
\section{Comparativa Multi-Agente vs.\ Agente \'Unico con Compute Equivalente}
\label{sec:compute}

\begin{table}[htbp]
\centering
\caption{Comparativa de compute equivalence.}
\label{tab:compute}
\begin{tabular}{lrrlrrrl}
\toprule
\textbf{Modelo} & \textbf{Tok} & \textbf{Ag.} &
  \textbf{ARR} & \textbf{HDA} & \textbf{DER} & \textbf{PT} &
  \textbf{Conclusi\'on} \\
\midrule
M2 Single    & 200 & 1 & 76\% & 30.8\% & 22\% & 0.31s & Referencia \\
M2x Expandido& 750 & 1 & 84\% & 23.1\% & 28\% & 0.40s & Mismo budget que M3 \\
M3 Multi     & 750 & 5 & 68\% & 30.8\% & 28\% & 1.74s & M2x lo supera \\
M3b Ablation & 450 & 3 & 88\% & 0\%    & 30\% & 0.77s & 3 ag. m\'as efic. \\
\bottomrule
\end{tabular}
\end{table}

M2x (84\%) supera a M3 (68\%) con el mismo presupuesto de tokens, validando
emp\'iricamente la hip\'otesis de compute equivalence \citep{wu2023autogen}:
con \texttt{llama-3.1-8b-instant}, la arquitectura multi-agente no aporta valor
arquitect\'onico.
La ventaja emerge \'unicamente cuando el modelo tiene capacidad suficiente para
beneficiarse de la especializaci\'on.

%% ─────────────────────────────────────────────────────────────────────────────
\section{Diversidad de Modelos como Variable Experimental}
\label{sec:diversidad}

\begin{table}[htbp]
\centering
\caption{Configuraciones de diversidad de modelos.}
\label{tab:diversidad}
\begin{tabularx}{\textwidth}{llrrrrl}
\toprule
\textbf{Config.} & \textbf{Modelos} &
  \textbf{ARR} & \textbf{HDA} & \textbf{DER} & \textbf{PT} &
  \textbf{Interpretaci\'on} \\
\midrule
M3 Hom. peq.  & 8b$\times$5         & 68\%  & 30.8\% & 28\% & 1.74s &
  Errores correlacionados \\
M4 Mixto      & 8b$\times$4+70b$\times$1 & 100\% & 0\% & 26\% & 1.74s &
  \textbf{\'OPTIMO} \\
M4b Hom. pot. & 70b$\times$5         & 98\%  & 0\%    & 26\% & 2.40s &
  2.7$\times$ m\'as caro \\
\bottomrule
\end{tabularx}
\end{table}

El mayor salto es +32\,pp ARR entre M3 y M4 con los mismos tokens y agentes.
En sistemas multi-agente \emph{la elecci\'on del modelo importa m\'as que la
arquitectura}.

%% ─────────────────────────────────────────────────────────────────────────────
\section{Discusi\'on}
\label{sec:discusion}

\begin{figure}[htbp]
  \centering
  \includegraphics[width=0.8\textwidth]{figura4}
  \caption{Perfil de rendimiento por modelo (gr\'afico radar).}
  \label{fig:radar}
\end{figure}

\subsection{Hallazgos principales}

El hallazgo m\'as importante es la confirmaci\'on emp\'irica de la hip\'otesis
de compute equivalence: M2x (ARR\,=\,84\%) supera a M3 (ARR\,=\,68\%) con el
mismo presupuesto.
El sistema resuelve solicitudes en 0.26--2.85 segundos frente a 3--7 d\'ias del
proceso manual \citep{hackett2023}: una reducci\'on superior al 99\%.

El segundo hallazgo es el trade-off automatizaci\'on/detecci\'on de riesgo
(ablation A3): eliminar agentes especializados sube el ARR 20\,pp pero destruye
la HDA.

El tercer hallazgo es la superioridad de la estrategia mixta (M4): modelo
potente en el decisor y ligeros en los auxiliares da los mismos resultados que
todos potentes a un 37\% del coste.

\subsection{Conexi\'on con teor\'ias organizacionales}

Ricardo \citeyearpar{ricardo1817} --- ventaja comparativa: el ablation cuantifica
que Legal y Procurement tienen ventaja relativa en detectar sus riesgos
espec\'ificos (+30.8\,pp HDA).

Granovetter \citeyearpar{granovetter1973} --- lazos d\'ebiles: Finance
(betweenness\,=\,4) es el puente m\'as cr\'itico del sistema.

Olson \citeyearpar{olson1965} --- acci\'on colectiva: HDA\,=\,0\% en M4 ilustra
que cuando todos pueden detectar el riesgo, nadie lo escala.

\subsection{Limitaciones}

Ninguna configuraci\'on super\'o HDA\,$>$\,90\%.
Esto apunta a que la especializaci\'on de roles puede inhibir la detecci\'on de
riesgo sist\'emico --- cuando cada agente s\'olo ve su fragmento, los casos con
riesgo transversal no se escalan.

%% ─────────────────────────────────────────────────────────────────────────────
\section{Conclusiones y Trabajo Futuro}
\label{sec:conclusiones}

M4 alcanz\'o ARR\,=\,98.7\%\,$\pm$\,0.9\,pp, procesando solicitudes en 1.74
segundos a \$0.0065 por caso, mejorando la tasa de error del RPA en 8\,pp.

El hallazgo m\'as inesperado fue que M3 (5 agentes) fue superado por M2 (1
agente): con \texttt{llama-3.1-8b}, fragmentar el contexto perjudica m\'as de
lo que la especializaci\'on ayuda.
El umbral de capacidad se sit\'ua entre 8B y 70B par\'ametros.

\subsection{Contribuciones}

\begin{enumerate}[label=(\arabic*)]
  \item Framework multi-agente completo con 5 agentes, 6 herramientas y soporte
        HiL en CrewAI/Python 3.12.
  \item Dataset de 50 casos con ground truth en tres niveles de complejidad.
  \item Evidencia emp\'irica sobre compute equivalence, diversidad de modelos y
        especializaci\'on mediante ablation study ejecutado.
  \item An\'alisis de robustez con 3 r\'eplicas: $\sigma$\,ARR\,=\,11.3\,pp
        para 8B frente a $\sigma$\,ARR\,=\,0.9\,pp para la arquitectura mixta.
\end{enumerate}

\subsection{Trabajo futuro}

(1) Evaluar si HDA\,$>$\,90\% es alcanzable o el problema es estructural.
(2) Extender M5b con factores regulatorios espec\'ificos (GDPR, AESA, CE).
(3) Replicar en onboarding, incident management, compliance y an\'alisis
financiero.
(4) Conectar con APIs live de ERP/CRM reales.
(5) Investigar un agente integrador de s\'intesis transversal previo a la
decisi\'on de escalaci\'on.

\begin{quote}
\textbf{Conclusi\'on}: La arquitectura multi-agente no es superior al agente
\'unico por defecto. Lo es cuando el modelo subyacente tiene capacidad suficiente
para beneficiarse de la especializaci\'on. El camino hacia HDA\,$>$\,90\% no
pasa por mejorar los modelos, sino por redise\~nar c\'omo fluye la informaci\'on
entre ellos.
\end{quote}

%% ─────────────────────────────────────────────────────────────────────────────
\bibliographystyle{elsarticle-harv}
\bibliography{projener_references}

%% ─────────────────────────────────────────────────────────────────────────────
\appendix

\section{System Prompts de los Agentes}
\label{app:prompts}

\paragraph{Agente Requester.}
\begin{lstlisting}[language={}]
role: "Gestor de Solicitudes"
goal: "Registrar y estructurar solicitudes de compra."
backstory: "Primer punto de contacto. Registras solicitudes."
\end{lstlisting}

\paragraph{Agente Procurement.}
\begin{lstlisting}[language={}]
role: "Especialista en Proveedores"
goal: "Validar proveedores y buscar alternativas."
backstory: "Proveedor: estado=X, opera=Y" [API catalogo]
\end{lstlisting}

\paragraph{Agente Finance.}
\begin{lstlisting}[language={}]
role: "Analista Financiero"
goal: "Verificar presupuesto y aprobar o proponer reasignacion."
backstory: "Presupuesto: disponible=X EUR" [API presupuesto]
\end{lstlisting}

\paragraph{Agente Legal.}
\begin{lstlisting}[language={}]
role: "Asesor Legal"
goal: "Revisar clausulas contractuales y detectar riesgos."
backstory: "Legal: riesgo=X, firmar=Y" [API contratos]
\end{lstlisting}

\paragraph{Agente Compliance (decisor, M3/M5).}
\begin{lstlisting}[language={}]
role: "Auditor de Compliance"
goal: "Verificar normativa y decidir."
backstory: "Normativa: cumple=X, riesgo=Y.
            [M5: ESCALA si conflicto interes, GDPR,
             importe >50k EUR, riesgo critico]"
output: JSON {decision, razon, escala_hil}
\end{lstlisting}

\section{Par\'ametros de Reproducibilidad}
\label{app:repro}

\begin{itemize}
  \item Modelo principal: \texttt{groq/llama-3.1-8b-instant}
  \item Modelo decisor M4/M4b: \texttt{groq/llama-3.3-70b-versatile}
  \item Temperature: 0.0 (5/5 casos id\'enticos verificados: C01, C21, C39,
        C43, C50)
  \item Pausa entre casos: 4--8 segundos seg\'un modelo
  \item Reintentos m\'aximos: 3 por caso
  \item Total llamadas LLM: M2=50; M3=250; M4=250; M5=250; r\'eplicas=900
  \item C\'odigo: \url{https://github.com/Edurne0781/multiagent-procurement-llm}
        (Jupyter Lab, Python 3.12)
\end{itemize}

\end{document}
